\section{מתודולוגיה}

\subsection{מודל מרקוב חבוי}

הרעיון של HMM הוא להניח שמאחורי הנתונים הנצפים קיימת סדרה של מצבים שאיננו רואים ישירות. במקרה שלנו איננו אומרים למודל מראש אילו ימים הם ``רגועים'' או ``לחוצים''. במקום זאת, המודל מנסה להסביר את וקטורי התכונות באמצעות מספר קטן של מצבים חבויים, כאשר לכל מצב יש התפלגות שונה של התצפיות ודפוס שונה של מעבר למצבים אחרים.

נסמן את המצב החבוי ביום \(t\) ב־\(S_t\), כאשר

\begin{equation}
S_t\in\{1,\ldots,K\}.
\end{equation}

הנחת מרקוב מסדר ראשון אומרת שלצורך חיזוי המצב הבא מספיק לדעת את המצב הנוכחי. אין צורך לזכור במפורש את כל רצף המצבים שקדמו לו:

\begin{equation}
P(S_t\mid S_1,\ldots,S_{t-1}) = P(S_t\mid S_{t-1}).
\end{equation}

את הסתברויות המעבר מרכזים במטריצה \(A\). כל איבר במטריצה הוא

\begin{equation}
a_{ij}=P(S_t=j\mid S_{t-1}=i).
\end{equation}

לדוגמה, אם \(a_{33}\) גבוה, המשמעות היא שכאשר המודל נמצא במצב 3 יש סיכוי גבוה שהוא יישאר במצב 3 גם ביום הבא. לכן הערכים שעל האלכסון של מטריצת המעבר מתארים את מידת ההתמדה של כל משטר.

ב־\tech{Gaussian HMM}, וקטור התכונות של כל יום מתואר באמצעות התפלגות גאוסית שתלויה במצב החבוי:

\begin{equation}
\mathbf{X}_t\mid S_t=k \sim \mathcal{N}(\boldsymbol{\mu}_k,\Sigma_k).
\end{equation}

כלומר, לכל מצב \(k\) יש וקטור ממוצעים משלו ומטריצת שונות־קווריאנס משלו. מצב שקט יכול למשל לקבל ממוצע תנודתיות נמוך, בעוד מצב לחוץ יכול לקבל תנודתיות וטווח יומי גבוהים יותר.

בפרויקט נעשה שימוש בספריית \tech{hmmlearn} \cite{hmmlearn} עם מטריצת קווריאנס מלאה. פרמטרי המודל נלמדים באמצעות \tech{Baum-Welch}, שהוא יישום של אלגוריתם \tech{Expectation-Maximization (EM)} \cite{baum1970maximization}. באופן אינטואיטיבי, האלגוריתם חוזר שוב ושוב על שני שלבים: תחילה הוא מעריך את ההסתברות שכל יום שייך לכל מצב; לאחר מכן הוא מעדכן את פרמטרי המצבים והמעברים בהתאם להסתברויות האלה. התהליך נמשך עד שהשיפור ב־\tech{log-likelihood} נעשה קטן והמודל מתכנס.

\subsection{בחירת מספר המצבים}

לפני שרצים על המספרים חשוב להבין מה מחפשים. מעט מדי מצבים עלולים לאחד יחד תקופות שונות מאוד, למשל תקופה רגועה ותקופה תנודתית. יותר מדי מצבים עלולים לפצל את הנתונים לקבוצות קטנות ולא יציבות שקשה לפרש. לכן מספר המצבים \(K\) הוא חלק מהמודל שיש לבחור באופן מבוקר.

לכל נכס נבדקו שלוש אפשרויות: \(K=2\), \(K=3\) ו־\(K=4\). לכל אפשרות בוצעו חמישה אתחולים אקראיים עם ה־\tech{seeds} הבאים: \(42,123,456,789,2026\). הסיבה לריבוי אתחולים היא שאלגוריתם EM יכול להתכנס לפתרונות מקומיים שונים בהתאם לנקודת ההתחלה.

קריטריון הבחירה הראשי היה \tech{Validation log-likelihood}. זהו מדד לשאלה עד כמה המודל שנלמד על \tech{Train} מסביר רצף חדש שלא שימש באימון. ערך גבוה יותר, כלומר פחות שלילי, נחשב טוב יותר. ה־\tech{Test} לא שימש כלל בבחירה.

בנוסף נשמרו \tech{AIC} ו־\tech{BIC} \cite{akaike1974new,schwarz1978estimating}. שני המדדים מאזנים בין איכות ההתאמה לבין מורכבות המודל: מודל עם יותר מצבים יכול להתאים טוב יותר לנתונים, אך הוא גם מכיל יותר פרמטרים ולכן מקבל קנס על מורכבות. ערך נמוך יותר של AIC או BIC עדיף. חשוב להדגיש שמדדים אלה בוחנים התאמה הסתברותית של המודל לנתונים; הם אינם מדדי חיזוי של כיוון היום הבא.

לבסוף נבדקו גם שכיחות השימוש בכל מצב, הנקראת \tech{occupancy}, והיציבות בין אתחולים. מצב שמופיע באחוז זעיר מאוד מהימים עלול להיות אמיתי, למשל מצב משברי נדיר, אך הוא גם עלול להעיד על פיצול יתר. לכן בחירת \(K\) אינה מסתמכת על מספר יחיד בלבד.

\subsection{מצב קשיח לעומת וקטור הסתברויות מלא}

לאחר אימון HMM יש שתי דרכים עיקריות לתאר את המצב של יום מסוים.

בגישה הקשיחה, \tech{hard state}, בוחרים רק את המצב בעל ההסתברות הגבוהה ביותר. לדוגמה, אם ההסתברויות לארבעה מצבים הן \(0.70,0.20,0.08,0.02\), היום יקבל את התווית ``מצב 1''. כל שאר המידע נזרק.

בגישה הרכה, \tech{soft posterior}, שומרים את כל וקטור ההסתברויות. בדוגמה הקודמת נשמר הווקטור המלא ולא רק המצב הראשון. היתרון הוא שהמודל יכול לבטא חוסר ודאות. יום שבו ההסתברויות הן \(0.52,0.45,0.02,0.01\) שונה מאוד מיום שבו הן \(0.98,0.01,0.01,0.00\), גם אם בשני המקרים המצב הקשיח זהה.

בהערכת ה־\tech{Test} השתמשנו ב־\textbf{הסקה סיבתית מהעבר בלבד}. עבור יום \(t\), המודל מקבל את התצפיות עד יום \(t\) ואינו רואה תצפיות עתידיות. וקטור ההסתברויות הוא לכן

\begin{equation}
\gamma_t(k)=P(S_t=k\mid X_{1:t}).
\end{equation}

הסימון \(X_{1:t}\) פירושו כל התצפיות מתחילת הרצף ועד היום הנוכחי. הוא אינו כולל את \(X_{t+1}\) או ימים מאוחרים יותר.

בגרסת התחזית הרכה, תחזית התשואה של היום הבא מחושבת כממוצע משוקלל על פני המצבים לפי \(\gamma_t(k)\). כך, יום שנמצא בגבול בין שני משטרים אינו מטופל כאילו המודל בטוח לחלוטין באחד מהם.

כדי למדוד את מידת אי־הוודאות של הווקטור כולו השתמשנו באנטרופיה מנורמלת:

\begin{equation}
H_t=-\frac{\sum_{k=1}^{K}\gamma_t(k)\log\gamma_t(k)}{\log K}.
\end{equation}

כאשר \(H_t\) קרוב לאפס, כמעט כל ההסתברות מרוכזת במצב אחד ולכן המודל בטוח יחסית. כאשר \(H_t\) גבוה, כמה מצבים מקבלים הסתברויות דומות ולכן המודל פחות בטוח בזהות המשטר.

\subsection{התמדה ומשך שהייה במצב}

אחד המאפיינים החשובים של משטר הוא לא רק איך הוא נראה, אלא גם כמה זמן הוא נוטה להימשך. ב־HMM מסדר ראשון ההסתברות לעזוב מצב מסוים קבועה בכל צעד, כל עוד נשארים באותו מצב. הנחה זו מובילה להתפלגות גאומטרית של משך השהייה.

במילים פשוטות, אם ההסתברות להישאר במצב \(i\) בכל יום היא \(a_{ii}\), אז בכל יום קיימת אותה הסתברות \(1-a_{ii}\) לצאת ממנו. מכאן שהמשך הממוצע הצפוי הוא

\begin{equation}
E[D_i]=\frac{1}{1-a_{ii}}.
\end{equation}

לדוגמה, אם \(a_{ii}=0.95\), משך השהייה הממוצע הצפוי הוא כ־20 ימים. בדוח משווים ערך תיאורטי זה למשך הממוצע שנמדד בפועל ברצף המצבים שפוענח.

ההשוואה הזו בודקת רק את הממוצע, לא את כל צורת התפלגות משכי השהייה. אם רוצים מודל שבו לכל מצב יש התפלגות משך גמישה יותר, אפשר בעתיד להשתמש ב־\tech{Hidden Semi-Markov Model} \cite{bulla2006application}.
